\section{Introduction}
\label{sec:introduction}
Protein structure predictors couple sequence information to geometry through
learned pairwise representations
\citep{jumper2021alphafold,gong2025protenix,seo2026atlasfold}.
An adapter can reuse a pretrained pair-generating operator, predicting
residue-level factor increments and combining them in its native coordinates.
We call this structured adapter the Factor head.
Freezing the predictor lets us study two inherited properties of this
construction: its restricted update form and its orientation relative to
downstream computation.

Parameter-efficient methods constrain weight updates, as in LoRA
\citep{hu2021lora}, or intervene on representations, as in ReFT
\citep{wu2024reft}. We ask: \emph{under a fixed task-training budget, does
output-channel rotation affect adaptation differently across parameterizations?}
Matching parameter counts and local spectra does not specify how a frozen
predictor uses an update. A comparison must also distinguish an orientation
effect within a head from its absolute quality and competitiveness with
another adapter.

Our four-cell design compares native and rotated versions of two heads.
Factor composes increments with native interface anchors through the frozen operator;
G+ receives the same extra sequence features and explicit anchors, with a free
affine output layer. Orthogonal transforms act only on completed residuals.
Each arm is trained separately under paired settings, beginning
at the same zero-residual function. For each head $H$, we measure the signed
Native--Rotated score difference $\Delta_H$ and compare heads through
$\Psi=\Delta_F-\Delta_{G+}$. A positive interaction means a larger signed
score difference for Factor; it establishes neither a larger absolute response
to rotation nor a positive Factor score difference by itself.

G+'s free output layer can absorb a rotation without changing its function
class, but this closure does not guarantee matching AdamW trajectories
\citep{loshchilov2019adamw}. Within-head rotations preserve local writer
spectra at identical inputs, interface anchors and parameters. Across heads,
we match the initial function,
approximate parameter count and training recipe, while Jacobians and
effective learning scales can differ. The interaction therefore compares
complete parameterizations without establishing an expressivity hierarchy.

Our fixed-budget studies establish positive interactions in several Protenix
recipes and in OpenFold \citep{ahdritz2024openfold} under the Train96/ESM2
recipe on observed panels. Fixed Protenix/ESMC models also establish the
primary interaction on 192 prospectively locked new targets (Fresh192).
The fixed OpenFold Train96 models do not establish their Fresh96 primary
interaction; OpenFold Train384 and AtlasFold likewise do not establish it
on their observed panels. Native Factor outperforms G+ in Protenix, whereas
G+ wins some OpenFold comparisons, so the interaction does not determine
which head performs best.

After Fresh192 was observed, a prespecified doubling of the Protenix update
budget retains a positive interaction while significantly reducing Factor's
own Native--Rotated score difference. The paired change in interaction remains
unestablished. A separate amended OpenFold schedule study on observed targets,
conditional on its fitted models and one fixed rotation, finds the largest
interaction with First-only injection, which also has the lowest absolute
quality. With All injection, a positive interaction coexists with a Factor
score difference whose interval includes zero and a negative G+ difference.
Thus evidence beyond the original stopping point supports the interaction,
while the schedule comparison demonstrates why maximizing $\Psi$ alone
cannot select an adapter.

Additional interventions test the limits of a mechanism account. We ask
whether adding shared orthogonal output freedom helps rotated Factor more
than Native Factor. The two preselected OpenFold rotated Factor arms improve
in both the original execution and a same-seed retraining repeat, but only
the original establishes an additional gain over Native's own improvement.
Across eight tested rotations, a prespecified teacher-update reconstruction
diagnostic did not establish the predicted ranking of post-training
Native--Rotated score differences.

Our contributions are the four-cell comparison and new-target test, evidence
that the interaction persists across two tested training budgets, and a
schedule comparison showing that interaction size does not select the best
absolute quality. The output-map and reconstruction studies test possible
explanations and delimit what the interaction identifies. Together, these
results motivate reporting absolute performance, both signed within-head
effects and their interaction under explicit training budgets.
